\scanpage{20}

Prouvons $(\mathrm{ii})\Longrightarrow(\mathrm{i})$. La question étant locale sur $X$, on peut supposer $X$ \uline{noethérien}.

On prouve par récurrence noethérienne que pour tte partie fermée $Y'$ de $Y$, $F$ a la propriété $P_1$ $(P_2)$ rel. à $Y'$, on peut supposer que pour tt $Y'$ fermé $\subset Y$, $Y'\ne Y$, $F$ a déjà la propriété en question \inserted{rel. à $Y'$} et se borner à prouver que $F$ a la propriété $P_1$ $(P_2)$ rel. à $Y$ dans ce cas.

\deleted{\unsure{Si|Le} $Y=\emptyset$, c'est trivial, sinon soit $\deleted{y}\inserted{x}$ un pt \uline{maximal} de $Y$.} On est ramené à \deleted{prouver que}, prouver que

\[
\Gamma(X,F)\longrightarrow\Gamma(\deleted{\unsure{X}}U,F)\qquad (U=X-Y)
\]

est injectif (bijectif). Soient $f,g\in\Gamma(X,F)$ ayant $=$ image dans $\Gamma(U,F)$, prouvons qu'ils sont identiques\unsure{;|,} \deleted{il suffit d'établir que \unsure{de se dire que} pour tt $x\in Y$, $f_x=g_x$ ; si $\overline x\ne Y$,}

% slanted
\margin{\deleted{\unsure{injectivité} \unsure{du|de} \illegible{1} \unsure{5}}}

\unsure{Cela.|Ceci.|local.} Tout d'abord, en vertu \deleted{\unsure{du lemme 5}} \inserted{et de l'hyp. $P_1$ \uline{en $x$}} ils sont identiques en $x$, donc au voisinage de $x$, et on achève par l'hyp. de récurrence.

Soit \unsure{Mainten\deleted{ant}|Maintenant} (\unsure{ds|sous|et} l'hyp $P_2$ en les $x\in Y$)

% archive sheet number: 19
